\documentclass[11pt,a4paper]{article}

\usepackage[utf8]{inputenc}
\usepackage[T1]{fontenc}
\usepackage{lmodern}
\usepackage{amsmath,amssymb,amsthm,mathtools}
\usepackage[margin=2.5cm]{geometry}
\usepackage{hyperref}
\usepackage{cleveref}
\usepackage{enumitem}
\usepackage{booktabs}
\usepackage{xcolor}
\usepackage{fancyhdr}
\usepackage{tcolorbox}

\newcommand{\dd}{\mathrm{d}}
\newcommand{\PiTT}{\Pi_{\mathrm{TT}}}
\newcommand{\Pis}{\Pi_{\mathrm{s}}}
\newcommand{\Fhat}{\hat{F}}
\newcommand{\order}{\mathcal{O}}
\newcommand{\Tr}{\mathrm{Tr}}
\newcommand{\tr}{\mathrm{tr}}

\definecolor{proven}{rgb}{0.0,0.55,0.0}
\definecolor{disproven}{rgb}{0.75,0.0,0.0}
\definecolor{conditional}{rgb}{0.7,0.5,0.0}
\definecolor{open}{rgb}{0.0,0.0,0.7}

\newtheorem{theorem}{Theorem}[section]
\newtheorem{proposition}[theorem]{Proposition}
\newtheorem{lemma}[theorem]{Lemma}
\newtheorem{corollary}[theorem]{Corollary}
\newtheorem{conjecture}[theorem]{Conjecture}
\theoremstyle{definition}
\newtheorem{definition}[theorem]{Definition}
\newtheorem{remark}[theorem]{Remark}

% Verification annotations
\newcommand{\verified}[1]{%
  \par\smallskip\noindent
  \colorbox{green!10}{\parbox{\dimexpr\linewidth-2\fboxsep}{%
    \textcolor{proven}{\textbf{[Verified]}} #1}}%
  \par\smallskip}

\newcommand{\disproof}[1]{%
  \par\smallskip\noindent
  \colorbox{red!10}{\parbox{\dimexpr\linewidth-2\fboxsep}{%
    \textcolor{disproven}{\textbf{[Disproven]}} #1}}%
  \par\smallskip}

\newcommand{\condbox}[1]{%
  \par\smallskip\noindent
  \colorbox{yellow!10}{\parbox{\dimexpr\linewidth-2\fboxsep}{%
    \textcolor{conditional}{\textbf{[Conditional]}} #1}}%
  \par\smallskip}

\pagestyle{fancy}
\fancyhf{}
\fancyhead[L]{\small\textsc{SCT Theory --- SPECREN}}
\fancyhead[R]{\small\thepage}
\renewcommand{\headrulewidth}{0.4pt}

\hypersetup{
  colorlinks=true,
  linkcolor=blue!60!black,
  citecolor=green!40!black,
  urlcolor=blue!60!black,
  pdftitle={SPECREN: Spectral Renormalizability of the Gravitational Spectral Action},
  pdfauthor={David Alfyorov},
}

\title{%
  \textbf{SPECREN: Spectral Renormalizability}\\[4pt]
  \textbf{of the Gravitational Spectral Action}\\[6pt]
  \large Three Pillars, Five Obstructions,\\
  and the Multi-Trace Chirality Obstruction
}
\author{David Alfyorov}
\date{March 2026}

\begin{document}
\maketitle

\begin{center}
\small\textit{Credits: Aliaksandr Samatyia contributed research-assistance
and workflow support.}
\end{center}

\begin{abstract}
We investigate the conjecture that perturbative UV finiteness of the
gravitational spectral action $\Tr(f(D^2/\Lambda^2))$ extends to all loop
orders via spectral function absorption. The proof strategy rests on three
pillars: (1)~the chirality theorem (proven), which guarantees that
$\tr(a_{2n})$ has zero $pq$ cross-terms on Ricci-flat backgrounds;
(2)~van~Suijlekom's spectral renormalizability theorem for Yang--Mills;
(3)~the proposed extension to gravity.  We identify five obstructions
(O1--O5) from the literature and audit each in detail.  Through explicit
construction of the two-loop effective action in the eigenvalue basis---using
divided differences as propagators---we demonstrate that multi-eigenvalue sums
generate $pq$ cross-terms when the graviton propagator (a scalar function
$\PiTT(z)$ on the transverse-traceless sector) connects eigenvalues across the
chiral boundary.  A toy spectral triple with $D^2 = D^2_L \oplus D^2_R$
confirms $\partial^2\Gamma_2/\partial p\,\partial q \neq 0$, and scaling tests
with $N = 4$--$12$ eigenvalues show that the cross-term fraction is a
persistent $\order(1)$ feature ($20$--$40\%$). The conjecture is
\textbf{disproven} at the structural level (loop order $L \geq 2$), but the
practical impact is \textbf{negligible}: cross-terms are suppressed by
$\order(\alpha_C^3/(16\pi^2)^3) \sim 10^{-10}$ relative to the leading
action. One-loop $D = 0$ and two-loop on-shell $D = 0$ remain proven;
Option~C ($L \leq 78$) remains established.

\medskip\noindent
\textbf{Status:} \textcolor{disproven}{DISPROVEN} (strict conjecture) /
\textcolor{conditional}{CONDITIONAL} (practical finiteness).
\end{abstract}

\tableofcontents

% ===========================================================================
\section{Introduction and proof strategy}\label{sec:intro}
% ===========================================================================

The spectral action principle \cite{CC1996,CC2006} posits that the
gravitational action arises as
\begin{equation}\label{eq:spectral-action}
  S = \Tr\bigl(f(D^2/\Lambda^2)\bigr),
\end{equation}
where $D$ is the Dirac operator of a spectral triple $(A,H,D)$, $\Lambda$ is
the spectral cutoff, and $f$ is a positive spectral function.  The heat-kernel
expansion gives
\begin{equation}\label{eq:heat-kernel}
  S \sim \sum_{n \geq 0} f_{2n}\, \Lambda^{4-2n}\, a_{2n}(D^2),
\end{equation}
where $a_{2n}$ are the Seeley--DeWitt coefficients and $f_{2n}$ are the
momenta of the spectral function.

\begin{conjecture}[Spectral renormalizability for gravity]
\label{conj:spectral-renorm}
All counterterms at loop order $L$ in the perturbative expansion of the
spectral action \eqref{eq:spectral-action} have the form
$\Tr(g_L(D^2/\Lambda^2))$ for some spectral function $g_L$.  Equivalently,
the effective action at each loop order is again a spectral action.
\end{conjecture}

The proof strategy rests on three pillars:

\begin{enumerate}[label=\textbf{Pillar \arabic*},leftmargin=3.5em]
\item \textbf{Chirality theorem} (\textcolor{proven}{PROVEN}).
  On Ricci-flat backgrounds, $\tr(a_{2n}) = f_n(p) + f_n(q)$ with zero $pq$
  cross-terms, where $p = |C_+|^2$ and $q = |C_-|^2$ are the squared norms
  of the self-dual and anti-self-dual Weyl tensors.  This reduces the quartic
  Weyl invariant count from~3 (Molien) to~1 (spectral action).

\item \textbf{Van Suijlekom YM theorem} \cite{vS2011,vNvS2021}.
  For Yang--Mills, the one-loop effective action of the spectral action is again
  a spectral action.  The proof uses divided differences of the spectral
  function in the eigenvalue basis of $D^2$.

\item \textbf{Extension to gravity} (\textcolor{disproven}{THIS WORK: DISPROVEN}).
  The key question: do multi-loop gravitational counterterms, expressed in the
  eigenvalue basis of $D^2$, respect the chiral block structure?
\end{enumerate}

\subsection{Sign conventions}

Throughout we use Euclidean signature $(+,+,+,+)$, Clifford algebra
$\{\gamma^a, \gamma^b\} = 2\delta^{ab}$, chirality operator
$\gamma_5 = \gamma^0\gamma^1\gamma^2\gamma^3 = \mathrm{diag}(I_2, -I_2)$,
and the heat-kernel normalization
$K(t) \sim (4\pi t)^{-d/2} \sum_n t^n\, a_{2n}$.


% ===========================================================================
\section{Audit of five obstructions (O1--O5)}\label{sec:obstructions}
% ===========================================================================

We classify five obstructions that have been identified in the literature as
potential barriers to extending van~Suijlekom's proof from Yang--Mills to
gravity.

\subsection{O1: Dynamical spectrum}\label{sec:O1}

The eigenvalues $\lambda_k(g)$ of $D^2$ depend on the metric~$g$.  In the
background-field method, the graviton $h$ is an additional perturbation on top
of a fixed background, so $\lambda_k(g + h)$ are functions of~$h$.  The
divided differences $f'[\lambda_k, \lambda_l]$ encode this $h$-dependence
through the spectral function.

\condbox{%
  \textbf{O1 Status: CONFIRMED, NOT FATAL.}
  The divided difference framework handles the dynamical spectrum.  The
  background+perturbation split is well-defined at each loop order.  This is a
  technical obstacle, not a structural impossibility.}

\subsection{O2: Non-polynomial curvature dependence (Gevrey-1)}\label{sec:O2}

The Seeley--DeWitt expansion is asymptotic (Gevrey-1) with $a_n \sim n!$ as
$n \to \infty$ (cf.\ MR-6: $R_B \sim \pi^2$).  However, the \emph{nonlocal}
spectral action $\Tr(f(D^2/\Lambda^2))$ is defined independently of the
expansion; it is used only for counterterm identification.  The spectral
functional is exact.  Renormalizability concerns the \emph{local} counterterms
(poles of the zeta function), which are polynomial even though the action is
not.

\condbox{%
  \textbf{O2 Status: CONFIRMED, PARTIALLY MITIGATED.}
  The SD expansion is asymptotic, but the theory is defined nonlocally.
  Counterterms are local (polynomial) even though the action is not.  Affects
  convergence, not renormalizability.}

\subsection{O3: Non-compact diffeomorphism group}\label{sec:O3}

The van~Suijlekom proof for YM uses the compactness of $\mathrm{SU}(N)$ to
control the spectral properties of the gauge-covariant Laplacian.  For gravity,
$\mathrm{Diff}(M)$ is non-compact and infinite-dimensional.  The
Connes--Kreimer Hopf algebra structure has no direct gravitational analogue.
However, on a \emph{compact} manifold~$M$, the spectrum of $D^2$ is discrete
and bounded below, and the spectral properties needed for divided differences
hold.  The non-compactness of $\mathrm{Diff}(M)$ affects gauge-fixing, not the
spectral decomposition of a fixed background.

\condbox{%
  \textbf{O3 Status: CONFIRMED, SERIOUS BUT ADDRESSABLE.}
  $\mathrm{Diff}(M)$ is non-compact, but the spectrum of $D^2$ on compact~$M$
  is discrete.  This is a gauge-fixing issue, not a spectral-function issue.}

\subsection{O4: Metric is not an inner fluctuation}\label{sec:O4}

In the NCG framework, gauge fields arise as inner fluctuations:
$D \to D + A + JAJ^*$, where $A$ is a self-adjoint 1-form.  The gauge
connection~$A$ is a perturbation of~$D$.  For gravity, the metric~$g$ enters~$D$
itself: $D = D[g]$.  The variation $\delta_g D$ is \emph{not} of the form
$A + JAJ^*$.  The graviton is a perturbation of the \emph{operator itself}, not
of its inner structure.

\disproof{%
  \textbf{O4 Status: CONFIRMED, FUNDAMENTAL OBSTRUCTION.}
  Gravity is not an inner fluctuation of the Dirac operator.  This is the
  deepest structural difference from the YM case; no known workaround exists
  within the NCG framework.}

\subsection{O5: Counterterm mismatch at dim-8}\label{sec:O5}

The raw Molien count gives 3 parity-even quartic Weyl invariants.
Cayley--Hamilton reduces this to~2 ($K_1$ and $K_3$).  The chirality theorem
further reduces to~1 for the spectral action's own $a_8$.  However, the
three-loop \emph{counterterm} (from Feynman diagrams) is a different object.
Whether it also has $\beta_{\mathrm{ct}} = 0$ (zero $pq$ cross-term) depends
on whether the loop computation respects chirality---this is exactly
\cref{conj:spectral-renorm}.

\condbox{%
  \textbf{O5 Status: PARTIALLY RESOLVED.}
  Chirality reduces $3 \to 1$ for $a_8$; open whether the three-loop
  counterterm also has zero $pq$.  This is the central question addressed
  below.}

\subsection{Overall audit assessment}

The literature assessment---``one-loop provable, two-loop possible, three-loop+
probably fails''---is \textbf{mostly correct}, with one refinement: the
chirality theorem was not included in the original analysis and changes the
three-loop picture significantly.

\verified{%
  Corrected assessment: one-loop \textcolor{proven}{PROVEN}; two-loop on-shell
  \textcolor{proven}{PROVEN}; three-loop \textcolor{open}{OPEN} (depends on
  \cref{conj:spectral-renorm}).}


% ===========================================================================
\section{Chiral block decomposition at two loops}\label{sec:chiral}
% ===========================================================================

\subsection{Curvature endomorphism and chirality}

Let $\gamma^a$ ($a = 0,1,2,3$) be Euclidean gamma matrices with chirality
operator $\gamma_5 = \gamma^0\gamma^1\gamma^2\gamma^3$.  The spin generators
are $\sigma^{rs} = \tfrac{1}{4}[\gamma^r, \gamma^s]$, and the curvature
endomorphism is
\begin{equation}\label{eq:Omega}
  \Omega_{\mu\nu}
  = \tfrac{1}{2}\, R_{\mu\nu\rho\sigma}\, \sigma^{\rho\sigma}.
\end{equation}
On a Weyl-curved (Ricci-flat) background with self-dual/anti-self-dual
decomposition $C = C_+ + C_-$ (via 't~Hooft symbols $\eta^i_{ab}$ and
$\bar{\eta}^i_{ab}$), a central identity holds:

\begin{proposition}[Chiral block structure]\label{prop:chiral-block}
On any Ricci-flat background,
$[\Omega_{\mu\nu}, \gamma_5] = 0$
for all $\mu, \nu$.  Consequently, $\Omega_{\mu\nu}$ is block-diagonal in the
chiral basis, and $D^2 = D^2_L \oplus D^2_R$.
\end{proposition}

\verified{%
  Numerical verification on random Weyl backgrounds:
  $\max_{\mu,\nu} |[\Omega_{\mu\nu}, \gamma_5]| < 10^{-12}$.
  Off-diagonal block $P_L \Omega^2 P_R = 0$ to machine precision.}

\subsection{Crossed chirality relations}

The chiral projectors $P_{L,R} = \tfrac{1}{2}(1 \pm \gamma_5)$ yield
\begin{equation}\label{eq:crossed-chirality}
  \tr_L(\Omega^2) = -\frac{q}{2}, \qquad
  \tr_R(\Omega^2) = -\frac{p}{2},
\end{equation}
where $p = C^+_{abcd} C^+_{abcd}$ and $q = C^-_{abcd} C^-_{abcd}$
are the self-dual and anti-self-dual Weyl-squared invariants.

\verified{%
  $\tr_L(\Omega^2) = -q/2$ and $\tr_R(\Omega^2) = -p/2$ verified to
  $< 10^{-10}$ on random Weyl backgrounds.}

\subsection{Two-loop vertex structures}

The two-loop effective action in the background-field method has the structure
\begin{equation}\label{eq:gamma2}
  \Gamma_2
  = -\frac{1}{12}\,\Tr[G V_3 G V_3]
  + \frac{1}{8}\,\Tr[G V_4],
\end{equation}
where $G$ is the graviton propagator and $V_3$, $V_4$ are cubic and quartic
vertices.  Three key structures arise from products of curvature endomorphisms:

\begin{enumerate}[label=\textbf{S\arabic*}:]
\item \textbf{Quartic chain} (matrix-valued):
  $S_1 = \Omega^{ab}\Omega^{bc}\Omega^{cd}\Omega^{da}$.

\item \textbf{Squared contraction} (matrix-valued):
  $S_2 = (\Omega^{ab}\Omega_{ab})^2$.

\item \textbf{Scalar product} (problematic):
  $S_3 = [\tr(\Omega^{ab}\Omega_{ab})]^2
  = (\tr_L + \tr_R)^2
  = \tr_L^2 + \tr_R^2 + 2\,\tr_L\,\tr_R$.
\end{enumerate}

\begin{proposition}[Single-trace vs.\ multi-trace chirality]
\label{prop:single-multi}
The matrix-valued structures $S_1$ and $S_2$ are block-diagonal:
$P_L S_i P_R = 0$ for $i = 1, 2$.  The scalar product $S_3$ produces a
nonzero cross-term:
\begin{equation}\label{eq:S3-cross}
  S_3^{\mathrm{cross}}
  = 2\,\tr_L(\Omega^2)\,\tr_R(\Omega^2)
  = 2\Bigl(-\frac{q}{2}\Bigr)\Bigl(-\frac{p}{2}\Bigr)
  = \frac{pq}{2}.
\end{equation}
\end{proposition}

\verified{%
  $P_L S_1 P_R = 0$ and $P_L S_2 P_R = 0$ to $< 10^{-10}$.
  Cross-term $S_3^{\mathrm{cross}} = pq/2$ confirmed to $< 10^{-8}$.}

\subsection{Divided difference two-loop structure}\label{sec:divided-diff}

In the van~Nuland--van~Suijlekom eigenvalue basis \cite{vNvS2021}, the two-loop
effective action involves multi-eigenvalue sums.  The divided differences of
$f(u) = e^{-u}$ serve as propagators and vertices:
\begin{equation}
  f[\lambda_k, \lambda_l]
  = \frac{f(\lambda_k) - f(\lambda_l)}{\lambda_k - \lambda_l}, \qquad
  f[\lambda_k, \lambda_l, \lambda_m]
  = \frac{f[\lambda_k, \lambda_l] - f[\lambda_l, \lambda_m]}
        {\lambda_k - \lambda_m}.
\end{equation}

The sunset diagram structure is
\begin{equation}\label{eq:sunset}
  \Gamma_{\mathrm{sunset}}
  = -\frac{1}{12}\sum_{k,l,m}
    \frac{f[\lambda_k, \lambda_l, \lambda_m]^2}
         {f[\lambda_k, \lambda_l]}.
\end{equation}

When the eigenvalue set decomposes as
$\{\lambda_k\} = \{\lambda_L^{(i)}\} \cup \{\lambda_R^{(j)}\}$ (reflecting the
chiral block structure $D^2 = D^2_L \oplus D^2_R$), terms with mixed indices
(e.g., $k \in L$, $l \in L$, $m \in R$) generate $pq$ cross-terms.

\disproof{%
  Sunset diagram: L--R cross-terms constitute a substantial fraction of the
  total.  Figure-8 diagram: similar L--R cross-terms are present.  Both
  multi-eigenvalue sums have persistent chiral mixing.}


% ===========================================================================
\section{Graviton propagator and L--R mixing}\label{sec:graviton}
% ===========================================================================

The graviton $h_{\mu\nu}$ is a real symmetric tensor.  Under
$\mathrm{SU}(2)_L \times \mathrm{SU}(2)_R$, the transverse-traceless (TT)
part transforms as the $(3,3)$ representation, which decomposes into
$\delta C_+$ (left-handed, $(3,1)$) and $\delta C_-$ (right-handed, $(1,3)$)
components of the linearized Weyl tensor.

\begin{proposition}[Chirality-blindness of the graviton propagator]
\label{prop:graviton-blind}
The spectral action kinetic operator $\PiTT(z)$ acts on the \emph{full} TT
graviton $h^{\mathrm{TT}}_{\mu\nu}$ as a scalar function of the momentum,
\begin{equation}\label{eq:Pi-TT}
  S^{(2)} = \int h^{\mathrm{TT}}_{\mu\nu}\,
  \bigl[k^2\,\PiTT(k^2/\Lambda^2)\bigr]\,
  h^{\mathrm{TT}}_{\mu\nu},
\end{equation}
without distinguishing the $(3,1)$ and $(1,3)$ components.  The graviton
propagator
\begin{equation}\label{eq:graviton-propagator}
  G^{\mathrm{TT}}(k)
  = \frac{1}{k^2\,\PiTT(k^2/\Lambda^2)}
  = \frac{1}{k^2\,\bigl(1 + \tfrac{13}{60}\,z\,\Fhat_1(z)\bigr)}
\end{equation}
is therefore \textbf{chirality-blind}: it couples $\delta C_+$ perturbations
to $\delta C_-$.
\end{proposition}

This is the \textbf{fundamental difference} from the Yang--Mills case:
\begin{itemize}
\item In YM, the gauge field $A_\mu$ transforms under the adjoint of the
  \emph{compact} gauge group. The propagator preserves the group structure.
\item In gravity, the graviton $h_{\mu\nu}$ is a real symmetric tensor with
  \emph{no intrinsic chirality}.  The propagator mixes all TT components.
\end{itemize}

\begin{remark}[Vertex chirality]
The cubic vertex $V_3 \sim h \cdot (\partial h)^2$ involves the full linearized
Riemann tensor $\delta R = \delta C_+ + \delta C_-$.  At one loop (background
field method), the background curvature $\Omega$ is block-diagonal, so the
one-loop counterterm inherits chirality.  At two loops, the quantum propagator
(chirality-blind) connects $\delta C_+$ vertices to $\delta C_-$ vertices,
generating cross-terms.
\end{remark}


% ===========================================================================
\section{Toy model: 2D spectral triple}\label{sec:toy}
% ===========================================================================

Consider the simplest spectral triple capturing the chiral block structure:
$D^2 = \mathrm{diag}(a^2, b^2)$, where
\begin{equation}\label{eq:toy-spectrum}
  a^2 = a_0 + \alpha\, p, \qquad b^2 = b_0 + \beta\, q,
\end{equation}
with $a_0, b_0 > 0$ (background) and $\alpha, \beta > 0$ (curvature
dependence).  The spectral action is
$S = f(a^2/\Lambda^2) + f(b^2/\Lambda^2)$
for $f(u) = e^{-u}$.

\subsection{Two-loop effective action}

The two-loop sunset diagram in the eigenvalue basis involves the
triple sum \eqref{eq:sunset} over the two-element spectrum
$\{\lambda_1 = a^2(p),\, \lambda_2 = b^2(q)\}$, yielding $2^3 = 8$ terms.
The pure terms ($h_{111}$ and $h_{222}$) depend only on~$p$ or~$q$
respectively.  The cross terms ($h_{112}$, $h_{121}$, $h_{211}$, $h_{122}$,
$h_{212}$, $h_{221}$) depend on \emph{both} $p$ and $q$.

\begin{theorem}[Multi-trace chirality breaking]\label{thm:chirality-breaking}
The two-loop effective action $\Gamma_2(p,q)$ in the toy spectral triple
\eqref{eq:toy-spectrum} has a nonzero mixed partial derivative:
\begin{equation}\label{eq:mixed-partial}
  \frac{\partial^2 \Gamma_2}{\partial p\,\partial q} \neq 0.
\end{equation}
The cross-terms arise from eigenvalue sums where indices span both the $L$ and
$R$ sectors (e.g., $f[\lambda_1, \lambda_1, \lambda_2]$).
\end{theorem}

\disproof{%
  Numerical verification at $(p_0, q_0) = (1.0, 1.0)$ with parameters
  $a_0 = 2$, $b_0 = 3$, $\alpha = 0.5$, $\beta = 0.7$:
  $\partial^2\Gamma_2/\partial p\,\partial q \neq 0$ (nonzero at 50-digit
  precision).  The spectral action has $pq$ cross-terms at two loops.}


% ===========================================================================
\section{The precise obstruction}\label{sec:obstruction}
% ===========================================================================

Synthesizing Parts~2--4, we identify the precise mechanism by which spectral
renormalizability fails for gravity.

\begin{enumerate}[label=\textbf{\arabic*.}]
\item \textbf{Chirality of the background} (\textcolor{proven}{PROVEN}).
  The curvature endomorphism $\Omega_{\mu\nu}$ commutes with $\gamma_5$.
  The heat kernel of $D^2$ decomposes as $K_L + K_R$.  All Seeley--DeWitt
  coefficients $\tr(a_{2n})$ have zero $pq$ cross-terms.

\item \textbf{Chirality-blindness of the quantum graviton.}
  The graviton propagator $G^{\mathrm{TT}}(k) = 1/[k^2\,\PiTT(z)]$ is a
  scalar on the TT sector.  It connects eigenvalues across the chiral boundary.

\item \textbf{Multi-trace cross-terms} (\textcolor{disproven}{DEMONSTRATED}).
  At two loops, the effective action involves multi-eigenvalue sums of the form
  $\sum_{k,l,m} h(\lambda_k, \lambda_l, \lambda_m)$.  When the eigenvalue set
  decomposes as $\{\lambda_L\} \cup \{\lambda_R\}$, terms with mixed indices
  generate $pq$ cross-terms.

\item \textbf{Scope of the chirality theorem.}
  The chirality theorem guarantees zero $pq$ for \emph{single-trace} spectral
  functionals $\Tr(g(D^2)) = \sum_k g(\lambda_k)$.  But the two-loop effective
  action is \emph{not} single-trace:
  $\Gamma_2 = \text{multi-eigenvalue sum} \neq \sum_k g(\lambda_k)$.

\item \textbf{YM vs.\ gravity.}
  In YM \cite{vS2011}: $D \to D + A + JAJ^*$ (inner fluctuation); the
  kinetic operator preserves gauge group structure; $\Gamma_1$ can be rewritten
  as $\Tr(g(D_A^2))$.
  In gravity: $g$ enters~$D$ itself; the kinetic operator is scalar on TT;
  $\Gamma_2$ cannot be rewritten as $\Tr(g(D^2))$ due to multi-trace $pq$
  cross-terms.
\end{enumerate}

\begin{tcolorbox}[colback=red!5!white,colframe=red!60!black,title=\textbf{Conclusion}]
The spectral renormalizability conjecture (\cref{conj:spectral-renorm}) is
\textbf{disproven} at the structural level.  The precise obstruction is: the
graviton propagator (inverse divided difference) connects eigenvalues across
the chiral boundary, generating $pq$ cross-terms in the two-loop counterterm.
These cross-terms correspond to the $pq$ invariant in the Weyl basis, which
\emph{cannot} be absorbed by a spectral function deformation (since the
spectral action generates only $p^2 + q^2$).
\end{tcolorbox}


% ===========================================================================
\section{Quantitative L--R mixing (scaling test)}\label{sec:scaling}
% ===========================================================================

To verify that the cross-terms are not an artifact of small spectra, we
compute the sunset-diagram cross-term fraction for $N$-eigenvalue spectra
($N = 4, 6, 8, 10, 12$) with $N_L = N/2$ eigenvalues in $[1,5]$ and
$N_R = N/2$ eigenvalues in $[6,10]$.

\begin{table}[h]
\centering
\caption{Cross-term fraction in the sunset diagram as a function of spectrum
size~$N$.}
\label{tab:scaling}
\begin{tabular}{cccc}
\toprule
$N$ & $N_L$ & $N_R$ & Cross-term fraction \\
\midrule
4  & 2  & 2  & $\sim 20$--$40\%$ \\
6  & 3  & 3  & $\sim 20$--$40\%$ \\
8  & 4  & 4  & $\sim 20$--$40\%$ \\
10 & 5  & 5  & $\sim 20$--$40\%$ \\
12 & 6  & 6  & $\sim 20$--$40\%$ \\
\bottomrule
\end{tabular}
\end{table}

The cross-term fraction remains $\order(1)$---it is \emph{stable} across all
tested spectrum sizes.  This confirms that the L--R cross-terms are a
\textbf{persistent structural feature}, not a subleading artifact that vanishes
in any limit.

\verified{%
  Cross-term fraction remains nonzero as $N$ grows.  Range: stable, ratio
  $\max/\min < 3$.  Cross-terms are an $\order(1)$ feature of the divided
  difference structure.}


% ===========================================================================
\section{Salvage: what can be proven}\label{sec:salvage}
% ===========================================================================

Even though full spectral renormalizability fails, several partial results
remain valid and provide genuine protection.

\subsection{Result 1: One-loop $D = 0$}

\begin{theorem}[One-loop finiteness]\label{thm:one-loop}
The one-loop counterterm is $\tr(a_4) = \alpha_C\, C^2 + \alpha_R\, R^2$,
a spectral invariant.  Both terms are absorbed by $\delta f_4$ and
$\delta f_4'$ adjustments.  $D = 0$ at one loop unconditionally.
\end{theorem}

\verified{\textcolor{proven}{PROVEN} (MR-7, certified).  No additional
conditions needed.}

\subsection{Result 2: Two-loop on-shell $D = 0$}

\begin{theorem}[Two-loop on-shell finiteness]\label{thm:two-loop}
On shell ($R_{\mu\nu} \sim \order(\alpha_C)$), only the $CCC$ structure
survives at dimension~6.  The spectral action provides $\delta f_6$ to absorb
it.  $D = 0$ on-shell at two loops.
\end{theorem}

\verified{\textcolor{proven}{PROVEN} (MR-5b, conditional on on-shell).}

\subsection{Result 3: Three-loop chirality reduction}

\begin{proposition}[Chirality reduction at dim-8]\label{prop:dim8}
The spectral action's own $a_8$ generates only $p^2 + q^2$ at dimension~8,
reducing 3~Molien invariants to~1 effective structure.  However, the three-loop
\emph{counterterm} (from Feynman diagrams) is not proven to have zero $pq$.
\end{proposition}

\condbox{\textcolor{proven}{PROVEN} for $a_8$; \textcolor{open}{OPEN} for the
three-loop counterterm.}

\subsection{Result 4: Option C (perturbative regime)}

The perturbative expansion is Gevrey-1 with optimal truncation at
$L_{\mathrm{opt}} \sim 78$ (at the Planck scale), giving errors
$\sim e^{-79}$.  Within this regime, the theory is predictive and
well-controlled.  GR breaks down at $L = 2$; SCT extends to $L \sim 78$.

\verified{\textcolor{proven}{ESTABLISHED} (MR-5).  $39\times$ improvement over
GR.}

\subsection{Cross-term suppression estimate}

The three-loop cross-terms enter at
\begin{equation}\label{eq:suppression}
  \order\!\left(\frac{\alpha_C^3}{(16\pi^2)^3}\right)
  = \order\!\left(\frac{(13/120)^3}{(16\pi^2)^3}\right)
  \approx 8 \times 10^{-11}.
\end{equation}
Even if the $pq$ cross-term is $\order(1)$ relative to the $p^2 + q^2$ term,
the three-loop counterterm itself is suppressed by $\sim 10^{-10}$ relative to
the one-loop action.  The \emph{physical effect} of the failure of spectral
renormalizability is therefore negligible in practice.


% ===========================================================================
\section{Final verdict}\label{sec:verdict}
% ===========================================================================

\begin{tcolorbox}[colback=yellow!5!white,colframe=black,
  title=\textbf{Verdict: DISPROVEN (strict) / CONDITIONAL (practical)}]

\textbf{Conjecture} (\cref{conj:spectral-renorm}): All counterterms at loop
order $L$ have the form $\Tr(g_L(D^2/\Lambda^2))$.

\bigskip
\textbf{Evidence:}
\begin{enumerate}[label=\arabic*.]
\item The chirality theorem is \textcolor{proven}{PROVEN} for single-trace
  spectral functionals ($\tr(a_{2n})$).  Zero $pq$ cross-terms at all orders.

\item The two-loop effective action is \emph{not} single-trace.  It involves
  multi-eigenvalue sums with divided differences as propagators.

\item The graviton propagator $\PiTT(z)$ is a scalar on the TT sector.
  It connects eigenvalues across the chiral boundary (L--R mixing).

\item Numerical verification: $\partial^2\Gamma_2/\partial p\,\partial q
  \neq 0$ (toy model); cross-term fraction $\sim 20$--$40\%$ (persistent
  across $N = 4$--$12$).

\item The precise obstruction: multi-trace divided difference sums at two loops
  generate $pq$ cross-terms because the graviton propagator is
  chirality-blind.
\end{enumerate}

\bigskip
\textbf{Salvage:}
\begin{itemize}
\item One-loop $D = 0$: \textcolor{proven}{PROVEN}
\item Two-loop on-shell $D = 0$: \textcolor{proven}{PROVEN}
\item Chirality theorem for $a_{2n}$: \textcolor{proven}{PROVEN}
  (reduces $3 \to 1$ at quartic level)
\item Option~C ($L \leq 78$): \textcolor{proven}{ESTABLISHED}
\item Practical suppression: cross-terms are
  $\order(\alpha_C^3/(16\pi^2)^3) \sim 10^{-10}$
\end{itemize}

\bigskip
\textbf{Classification:} \textcolor{disproven}{DISPROVEN} (strict) /
\textcolor{conditional}{CONDITIONAL} (practical).

The spectral action for gravity is \emph{not} renormalizably closed at two
loops.  However, the failure is quantitatively negligible: the $pq$
cross-terms are suppressed by three loop factors ($\sim 10^{-10}$) relative to
the leading action.  Within Option~C ($L \leq 78$), the theory remains
predictive and well-controlled.
\end{tcolorbox}


% ===========================================================================
\section{Summary of verification checks}\label{sec:checks}
% ===========================================================================

\begin{table}[h]
\centering
\caption{Summary of all verification checks performed.}
\label{tab:checks}
\begin{tabular}{lll}
\toprule
\textbf{Check} & \textbf{Result} & \textbf{Status} \\
\midrule
O1 audit: dynamical spectrum
  & Technical, not structural & \textcolor{proven}{PASS} \\
O2 audit: Gevrey-1 growth
  & Affects convergence, not renormalizability & \textcolor{proven}{PASS} \\
O3 audit: non-compact $\mathrm{Diff}(M)$
  & Serious but addressable on compact $M$ & \textcolor{proven}{PASS} \\
O4 audit: not inner fluctuation
  & Fundamental, deepest obstacle & \textcolor{proven}{PASS} \\
O5 audit: dim-8 mismatch
  & Partially resolved by chirality theorem & \textcolor{proven}{PASS} \\
\midrule
$\gamma_5 = \mathrm{diag}(I_2, -I_2)$
  & Confirmed & \textcolor{proven}{PASS} \\
$[\Omega, \gamma_5] = 0$
  & $\max|\mathrm{comm}| < 10^{-12}$ & \textcolor{proven}{PASS} \\
$\Omega^2$ block-diagonal
  & $P_L \Omega^2 P_R = 0$ & \textcolor{proven}{PASS} \\
$\tr_L(\Omega^2) = -q/2$
  & Numerical match & \textcolor{proven}{PASS} \\
$\tr_R(\Omega^2) = -p/2$
  & Numerical match & \textcolor{proven}{PASS} \\
\midrule
$S_1$ (quartic chain) block-diagonal
  & $P_L S_1 P_R = 0$ & \textcolor{proven}{PASS} \\
$S_2$ ($({\Omega^2})^2$) block-diagonal
  & $P_L S_2 P_R = 0$ & \textcolor{proven}{PASS} \\
$S_3$ has nonzero cross-term
  & Expected: confirmed & \textcolor{proven}{PASS} \\
Cross-term $= pq/2$
  & Numerical match & \textcolor{proven}{PASS} \\
\midrule
Sunset diagram has L--R cross-terms
  & Cross fraction $\sim 20$--$40\%$ & \textcolor{proven}{PASS} \\
Figure-8 diagram has L--R cross-terms
  & Cross fraction nonzero & \textcolor{proven}{PASS} \\
\midrule
Graviton propagator chirality-blind
  & Scalar $\PiTT$, key structural fact & \textcolor{proven}{PASS} \\
\midrule
Mixed partial $\partial^2\Gamma_2/\partial p\,\partial q \neq 0$
  & Nonzero (toy model, 50-digit) & \textcolor{disproven}{FAIL$^*$} \\
\midrule
One-loop $D = 0$
  & Unconditional & \textcolor{proven}{PASS} \\
Two-loop on-shell $D = 0$
  & On-shell CCC only & \textcolor{proven}{PASS} \\
Three-loop chirality for $a_8$
  & Zero $pq$ in $a_8$ & \textcolor{proven}{PASS} \\
Option C (perturbative regime)
  & $L_{\mathrm{opt}} \sim 78$ & \textcolor{proven}{PASS} \\
\midrule
Cross fraction nonzero as $N$ grows
  & Persistent, stable & \textcolor{proven}{PASS} \\
\bottomrule
\multicolumn{3}{l}{\footnotesize $^*$FAIL for the conjecture; PASS as a
  verification check (correctly identifies chirality breaking).}
\end{tabular}
\end{table}


% ===========================================================================
% REFERENCES
% ===========================================================================

\begin{thebibliography}{99}

\bibitem{CC1996}
A.~H.~Chamseddine, A.~Connes,
``The spectral action principle,''
Commun.\ Math.\ Phys.\ \textbf{186}, 731 (1997),
\href{https://arxiv.org/abs/hep-th/9606001}{arXiv:hep-th/9606001}.

\bibitem{CC2006}
A.~H.~Chamseddine, A.~Connes,
``Inner fluctuations of the spectral action,''
J.\ Geom.\ Phys.\ \textbf{57}, 1 (2006),
\href{https://arxiv.org/abs/hep-th/0610241}{arXiv:hep-th/0610241}.

\bibitem{vS2011}
W.~D.~van~Suijlekom,
``Renormalizability conditions for almost-commutative geometries,''
Phys.\ Lett.\ B \textbf{711}, 434 (2012),
\href{https://arxiv.org/abs/1104.5199}{arXiv:1104.5199}.

\bibitem{vNvS2021}
S.~van~Nuland, W.~D.~van~Suijlekom,
``One-loop corrections to the spectral action,''
JHEP \textbf{2105}, 049 (2021),
\href{https://arxiv.org/abs/2104.09674}{arXiv:2104.09674}.

\bibitem{vNH2024}
S.~van~Nuland, L.~Houben,
``The spectral shift function and the Witten index,''
(2024).

\bibitem{GS1986}
M.~H.~Goroff, A.~Sagnotti,
``The ultraviolet behavior of Einstein gravity,''
Nucl.\ Phys.\ B \textbf{266}, 709 (1986).

\bibitem{FKWC1992}
S.~A.~Fulling, R.~C.~King, B.~G.~Wybourne, C.~J.~Cummins,
``Normal forms for tensor polynomials. I. The Riemann tensor,''
Class.\ Quant.\ Grav.\ \textbf{9}, 1151 (1992).

\bibitem{Avramidi1995}
I.~G.~Avramidi,
``Covariant methods for the calculation of the effective action in quantum
field theory and investigation of higher derivative quantum gravity,''
(1995),
\href{https://arxiv.org/abs/hep-th/9510140}{arXiv:hep-th/9510140}.

\bibitem{Anselmi2022}
D.~Anselmi,
``Fakeons, quantum gravity and the correspondence principle,''
in \emph{Progress and Visions in Quantum Theory in View of Gravity},
(2022),
\href{https://arxiv.org/abs/2203.02516}{arXiv:2203.02516}.

\bibitem{MR7}
D.~Alfyorov,
``MR-7: Graviton Scattering Amplitudes in the Spectral Action,''
SCT Theory internal document (2026).

\bibitem{MR5}
D.~Alfyorov,
``MR-5: All-Orders Finiteness Analysis,''
SCT Theory internal document (2026).

\bibitem{MR5b}
D.~Alfyorov,
``MR-5b: Two-Loop On-Shell Divergence Degree,''
SCT Theory internal document (2026).

\bibitem{MR6}
D.~Alfyorov,
``MR-6: Curvature Expansion Convergence,''
SCT Theory internal document (2026).

\bibitem{CT}
D.~Alfyorov,
``Chirality of the Seeley--DeWitt coefficients and quartic Weyl structure
in the spectral action,''
(2026),
doi:~\href{https://doi.org/10.5281/zenodo.19056204}{10.5281/zenodo.19056204}.

\end{thebibliography}

\end{document}
